
\textbf{Proposition 4.3.1.} \uline{Avec les notations ci-dessus, pour toute suite} @@M0@@ \uline{d'entiers positifs avec} @@M1@@\uline{, le coefficient de}
@@M2@@
\uline{dans le polynôme de Snapper} @@M3@@ \uline{est égal à}
@@M4@@
\uline{où} @@M5@@ \uline{est la classe de} @@M6@@ \uline{dans} @@M7@@\uline{. Donc dans le développement ordinaire de ce polynôme suivant les monômes en des} @@M8@@ @@M9@@\uline{, le coefficient du monôme} @@M10@@ \uline{est égal à}
@@M11@@

En particulier on obtient immédiatement le

\textbf{Corollaire 4.3.3.} \uline{Le coefficient du monôme} @@M12@@ \uline{dans le polynôme de Snapper} @@M13@@ \uline{est égal à}
@@M14@@

\emph{Démonstration} : Posons @@M15@@ pour tout @@M16@@. Tenant compte des définitions on trouve immédiatement que l'entier (4.3.2) est égal à @@M17@@. Lorsqu'on substitue dans (2.1.4) @@M18@@, @@M19@@, alors (2.1.5) donne, (compte tenu de ce que @@M20@@ pour tout @@M21@@) la congruence
@@M22@@
Donc le coefficient de @@M23@@ dans le polynôme @@M24@@ est égal à @@M25@@, C.Q.F.D.

